% ============================================
% LUC-14
% Extraído de main.pandoc.tex
% ============================================
\fragtitle{luc}{L}{LUC}{14}{Autoavaliação Clínica (DSM-5-TR)}
\markboth{LUC-14 --- Autoavaliação Clínica (DSM-5-TR)}{}

\noindent \textit{Documento extraído do arquivo pericial.}

\vspace{0.5em}\noindent\rule{\textwidth}{0.5pt}\vspace{0.5em}

\noindent {\footnotesize\texttt{ARQUIVO P\'UBLICO $\cdot$ PROVA DOCUMENTAL $\cdot$ PROCESSO 1.260}}\\[-0.2em]

\noindent\textbf{Documento:} Autoavaliação diagnóstica não solicitada.\\
\noindent \textbf{Elaborado por:} Dra. Lúcia Mendes — CRM-MG 28.391 — Psiquiatra Forense.\\
\noindent\textbf{Data:} 13 de abril de 2026 (05:47 -- 06:31).\\
\noindent \textbf{Solicitante:} a própria avaliada.\\
\noindent \textbf{Finalidade:} verificar correspondência entre sintomas atuais e critérios
\noindent diagnósticos do DSM-5-TR para Episódio Psicótico Breve (298.8),
\noindent Transtorno Dissociativo de Identidade (300.14) e Transtorno Psicótico
\noindent Compartilhado (Folie à Deux — 297.3).\\
\noindent \textbf{Status do documento:} não finalizado. A avaliada interrompeu o
\noindent preenchimento na página 4.

\noindent \textbf{CRITÉRIOS DIAGNÓSTICOS --- EPISÓDIO PSICÓTICO BREVE (DSM-5-TR 298.8)}

\textbf{A. Presença de um (ou mais) dos seguintes sintomas.}
\noindent Pelo menos um deve ser (1), (2) ou (3).

\noindent \textbf{(1) Delírios.}

\noindent Acredito que uma entidade\slash inteligência\slash padrão
\noindent contaminou um paciente em 1917, propagou-se por seu diário,
\noindent infectou um médico em 1979 e agora me contaminou pela leitura.

\noindent Reconheço, simultaneamente, que esta crença satisfaz
\noindent todos os critérios para ideação delirante: é fixa, idiossincrática,
\noindent impermeável a contra-evidência e causa sofrimento clinicamente significativo.

\noindent \textbf{Critério A(1): presente.}

\noindent \textbf{(2) Alucinações.}

\noindent Relato alucinação olfativa persistente (odor doce,
\noindent adocicado, descrito como ``flor murcha'' ou ``carne passada'').

\noindent Relato alucinação visual intermitente (percepção de
\noindent brilho amarelo no canto do campo visual, sem correspondente físico).

\noindent Relato alucinação tátil (sensação de peso no tórax,
\noindent independente de posição corporal, sem causa cardíaca ou respiratória
\noindent identificável).

\noindent Relato alteração na temperatura corporal localizada
\noindent (sensação de calor na região ocular direita, sem hipertermia objetiva).

\noindent Relato alteração proprioceptiva (sensação de ``mão
\noindent que se move sozinha'' durante a escrita manual, sem registro
\noindent eletromiográfico).

\noindent \textbf{Nota clínica:} todas as modalidades alucinatórias
\noindent relatadas correspondem exatamente às descritas pelo paciente índice
\noindent (Joaquim Antônio Correia, 1.260) e pelo caso secundário (Dr. Heitor
\noindent Oliveira). A coincidência interindividual de conteúdo, ordem de
\noindent aparecimento e progressão temporal inviabiliza a hipótese de
\noindent sugestão independente.

\noindent Considero esta nota clínica uma tentativa desesperada de
\noindent transformar observação em explicação.

\noindent \textbf{Critério A(2): presente.}

\noindent \textbf{(3) Discurso desorganizado.}

\noindent Este documento é a prova. Uma profissional com 15 anos
\noindent \textit{de experiência escrevendo uma autoavaliação diagnóstica }para si mesma
\noindent às 5:47 da manhã, em casa, sem supervisão, sem solicitação, sem
\noindent finalidade legal. Se uma paciente minha me entregasse este documento,
\noindent eu anotaria no prontuário: ``discurso superelaborado com conteúdo
\noindent persecutório, insight parcial, risco de deterioração.''

\noindent \textbf{Critério A(3): presente.}

\noindent \textbf{(4) Comportamento grosseiramente desorganizado.}

\noindent Não aplicável. A avaliada manteve funcionamento
\noindent ocupacional nos últimos 7 dias. Compareceu a compromissos. Atendeu
\noindent telefonemas. Respondeu e-mails. Nenhum dos interlocutores relatou
\noindent comportamento incomum.

\noindent \textbf{Critério A(4): ausente.}

\noindent \textbf{B. Duração do episódio: pelo menos um dia e menos de um mês.}

\noindent Início dos sintomas: 5 de abril de 2026 (primeira leitura
\noindent do diário). Data atual: 13 de abril de 2026. Duração: 8 dias.

\noindent \textbf{Critério B: presente.}

\noindent \textbf{C. Exclusão de transtorno afetivo.}

\noindent Sem histórico de episódios maníacos ou depressivos
\noindent maiores em 15 anos de prática clínica e 42 anos de vida. Sem
\noindent antecedentes familiares de primeiro grau.

\noindent \textbf{Critério C: presente.}

\noindent \textbf{D. Exclusão de substância/condição médica.}

\noindent Sem uso de substâncias psicoativas. Sem medicação em
\noindent curso. Exames laboratoriais solicitados em 09/04 (hemograma,
\noindent função tireoidiana, vitamina B12) dentro da normalidade.

\noindent \textbf{Critério D: presente.}

\noindent \textbf{E. Exclusão de outra condição psiquiátrica prévia.}

\noindent Sem diagnóstico prévio de transtorno psicótico,
\noindent dissociativo ou de personalidade. Histórico de ansiedade generalizada
\noindent leve (2018--2020), remitida sem medicação.

\noindent \textbf{Critério E: presente.}

\noindent \textbf{DIAGNÓSTICO DIFERENCIAL}

\noindent \textbf{Hipótese 1: Episódio Psicótico Breve (298.8).}

\noindent Todos os critérios preenchidos. Seria o diagnóstico
\noindent mais parcimonioso --- um episódio reativo ao contato com material
\noindent narrativo de alto impacto traumático, com conteúdo delirante
\noindent moldado pelo material de exposição.

\noindent Contraindicação: a coincidência interindividual de
\noindent sintomas entre três pessoas (Joaquim, Oliveira, eu) que não tiveram
\noindent contato entre si. Episódios Psicóticos Breves são individualmente
\noindent determinados. Não produzem réplicas sintomáticas em série.

\textbf{Hipótese 2: Transtorno Psicótico Compartilhado}

\noindent Tradicionalmente, exige contato direto entre o
\noindent paciente índice e o(s) secundário(s). Joaquim morreu em 1979.
\noindent Oliveira morreu em algum momento entre 1979 e 2026. Eu nunca
\noindent conheci nenhum dos dois. A transmissão se deu exclusivamente
\noindent pelo texto.

\noindent Não há precedente no DSM para folie à deux mediada
\noindent por leitura. Mas também não há exclusão explícita.

\noindent Se eu estiver certa, o diário é o primeiro caso
\noindent documentado de transtorno psicótico transmitido por narrativa.

\noindent Se eu estiver errada, o diário é um gatilho narrativo
\noindent poderoso o bastante para produzir o mesmo surto em três pessoas
\noindent sequencialmente. Uma coincidência que 15 anos de estatística clínica
\noindent me ensinaram a não aceitar.

\noindent \textbf{Hipótese 3: A entidade existe.}

\noindent Não é um diagnóstico. É uma constatação que não
\noindent cabe em CID nem em DSM. Mas a observação clínica, quando bem
\noindent feita, não descarta o que não consegue classificar. Apenas registra.

\noindent Estou registrando.

\noindent \textbf{PLANO TERAPÊUTICO}

\noindent 1. Devolver o diário ao arquivo. Não abrir novamente.

\noindent 2. Solicitar avaliação psiquiátrica independente.

\noindent 3. Iniciar antipsicótico atípico se os sintomas
\noindent persistirem por mais de 14 dias.

\noindent 4. Suspender leitura de qualquer material relacionado
\noindent ao caso.

\noindent 5. Não escrever mais nada à mão.

\noindent 6. Descartar este documento antes que alguém encontre.

\noindent 7. Não contar a ninguém o que está acontecendo.

\noindent \textbf{EVOLUÇÃO (06:31)}

\noindent Interrompi o preenchimento há 23 minutos. Reli o que escrevi.
\noindent Identifiquei três erros.

\noindent O primeiro erro: a redação está na primeira voz de Joaquim antes do
\noindent internamento. A estrutura frasal --- tópico + justificativa + ponderação ---
\noindent é a mesma que ele usa nas primeiras páginas do diário.

\noindent Estou escrevendo como ele escrevia.

\noindent O segundo erro: o plano terapêutico está em contradição com a
\noindent hipótese 3. Se a entidade existe, devolver o diário não interrompe
\noindent o contágio. O contágio já está em mim. O diário é apenas o vetor.

\noindent O vetor já cumpriu sua função.

\noindent \textit{O terceiro erro: estou escrevendo este documento contra a orientação do meu próprio plano terapêutico.}

\noindent O plano diz: ``Não escrever mais nada à mão.'' No momento em que escrevi a linha 7 do plano, já estava violando a linha 5.

\noindent A mão escreveu antes que eu decidisse.

\noindent A mão \textit{sabia} que eu ia continuar.

\noindent \textit{O arquivo \texttt{L260-14A.pdf} foi encontrado no computador da Dra. Lúcia
\noindent Mendes em 14 de maio de 2026, durante a varredura forense do equipamento. O
\noindent documento estava salvo na área de trabalho, sem criptografia, com o nome
\texttt{nao\_apagar\_isto.docx}. A análise forense confirma que o arquivo foi criado
\noindent em 13 de abril e modificado pela última vez em 12 de maio, às 2:47, quando
\noindent Lúcia anexou o e-mail abaixo. Não há outro registro de acesso posterior. O
\noindent documento não foi enviado a ninguém até 12 de maio. A avaliada não seguiu
\noindent nenhuma das recomendações do próprio plano terapêutico.}

\noindent \textbf{Maio de 2026 --- Belo Horizonte}

\noindent O último e-mail de Dra. Lúcia Mendes foi enviado em 12 de maio de 2026, às 2:47 da manhã.
\noindent O destinatário era o Dr. Ronaldo Torres, do Arquivo Público Mineiro.
\noindent \textquotedblleft{\textit{Ronaldo,}}
\noindent \textit{Preciso que você guarde este relato. Se algo acontecer comigo, quero que fique registrado.}
\noindent \textit{Li o diário completo do paciente 1.260. São 62 anos de registro --- de 1917 a 1979. As últimas páginas --- segundo o prontuário e o relato do coveiro --- teriam sido escritas depois da morte dele. Eu sei que é impossível, e é por isso que não coloco uma frase dessas em nenhum laudo oficial. Mas as páginas estão lá. A datação da tinta é inconclusiva. A caligrafia é a mesma.}
\noindent \textit{Joaquim Antônio Correia não estava louco.}
\noindent \textit{Ele estava possuído.}
\noindent \textit{Eu sei que isso soa absurdo vindo de uma psiquiatra forense. Eu passei 15 anos diagnosticando traumas, dissociações, psicoses. Eu acreditava que tudo tinha uma explicação neurológica ou psicológica.}
\noindent \textit{Agora acredito que existem coisas que a psicologia não explica.}
\noindent \textit{Li o diário. Três dias depois, comecei a sonhar com a vala. Uma semana depois, senti o cheiro. Duas semanas depois, meu olho direito começou a clarear.}
\noindent \textit{O mesmo olho. O mesmo amarelo.}
\noindent \textit{Fiz exames. Tomografia. Ressonância. Sangue. Não há explicação médica.}
\noindent \textit{Fui ao Colônia. Encontrei um caderno de Joaquim escondido numa cela --- dos anos de mutismo. Encontrei um velho que trabalhou lá nos anos 50. Ele me disse: 'A senhora leu o diário. Agora é sua vez.'}
\noindent \textit{Perguntei a vez do quê. Ele disse: 'De hospedar.'}
\noindent \textit{Estou com medo, Ronaldo. O medo não passa pela razão. Fica no corpo. Um medo que vem do estômago. Um medo que tem cheiro doce.}
\noindent \textit{Hoje, olhei no espelho e vi o olho amarelo me olhando de volta.}
\noindent \textit{Não era o meu olho.}
\noindent \textit{Era o olho dele.}
\noindent \textit{Olho de Joaquim.}
\noindent \textit{Ou daquilo que estava dentro dele.}
\noindent \textit{Vou entregar o diário no arquivo amanhã. Vou colocar num pano preto e deixar na recepção. Depois disso, vou sumir por um tempo.}
\noindent \textit{Não procure por mim.}
\noindent \textit{Se alguém perguntar, diga que o diário foi devolvido.}
\noindent \textit{E que o ciclo continua.}
\noindent Lúcia\textquotedblright{}
\noindent O diário foi encontrado na recepção do Arquivo Público Mineiro em 13 de maio de 2026, envolto num pano preto.
\noindent Dra. Lúcia Mendes não foi vista desde então.
\noindent Seu apartamento foi encontrado vazio. A cama estava arrumada. O café estava frio na xícara. O computador estava aberto na página de pesquisa: \textquotedblleft{\textit{Como passar o diário adiante sem ler?}\textquotedblright}
\noindent Não havia resposta.
\noindent Não há resposta.
\noindent O ciclo não tem reversa.
\noindent A única saída é entrar.
\noindent \textit{Nota do Arquivo Público Mineiro:}
\noindent \textit{O diário do Paciente 1.260 encontra-se atualmente em nossa custódia.}
\noindent \textit{Aguardando o próximo leitor.}
\noindent \textit{Aguardando você.}
\noindent \textbf{Abril de 2026 --- Belo Horizonte / Quixeramobim}

\noindent Depois de ler a carta de Oliveira, Lúcia fez o que qualquer psiquiatra forense faria: seguiu o rastro.

\noindent O Dr. Heitor Oliveira desapareceu em 13 de junho de 1979. Seu carro foi encontrado três semanas depois num estacionamento em Quixeramobim, Ceará.

\noindent A gasolina tinha acabado. As portas estavam destrancadas. Dentro, uma mochila com roupas, um caderno com 32 páginas escritas (que ninguém leu), e uma passagem de ônibus para Barbacena com data de 17 de junho --- o mesmo dia em que ele chegou.

\noindent Ele tinha voltado.
\noindent Lúcia passou uma tarde no Arquivo Público, revirando caixas. Encontrou o caderno de Oliveira --- esquecido, empoeirado, com a etiqueta \textquotedblleft{\textit{Pertences do Dr. H. Oliveira --- 1979\textquotedblright}}

\noindent Ela leu as 32 páginas em 40 minutos.

\noindent Eram anotações clínicas, mas com uma diferença: a letra mudava ao longo das páginas. Começava firme, médica. Terminava trêmula, irregular, \textit{infantil}.
\noindent A última frase era:
\noindent \textquotedblleft{\textit{Voltei para Barbacena. O ciclo não me deixa ir. Devolvo o diário ao arquivo. Espero que a próxima leitora entenda melhor do que eu.\textquotedblright{}}}

\noindent \textbf{A testemunha.}
\noindent No dia seguinte, Lúcia encontrou Dona Maria Cândida --- 92 anos, ex-enfermeira do Colônia, hoje residente em Barbacena.
\noindent A entrevista durou duas horas.
\noindent \textquotedblleft{\textit{O senhor Oliveira? Lembro. Ele chegou em 79, pouco antes do paciente 1.260 morrer. Era diferente dos outros médicos. Fazia perguntas. Olhava nos olhos. Os outros...\textquotedblright{}} Ela fez uma pausa. }\textquotedblleft{\textit{Os outros só aplicavam choque.\textquotedblright{}}}
\noindent \textquotedblleft{\textit{E depois da morte do Joaquim?\textquotedblright{}}}
\noindent \textquotedblleft{\textit{Ele mudou. Ficou estranho. Passava horas na ala dos silenciados, mesmo sem nenhum paciente lá. Uma vez eu vi ele falando sozinho na cela vazia.\textquotedblright{}}}
\noindent \textquotedblleft{\textit{O que ele dizia?\textquotedblright{}}}
\noindent Dona Maria Cândida hesitou.
\noindent \textquotedblleft{\textit{Dizia: 'Eu sei que a senhora está aqui. Eu sei que a senhora me ouve. Eu sei que a senhora não é uma história.'\textquotedblright{}}}
\noindent Lúcia anotou.
\noindent \textquotedblleft{\textit{Mais alguma coisa?\textquotedblright{}}}
\noindent \textquotedblleft{\textit{No dia em que ele sumiu, ele deixou um bilhete na minha mesa.\textquotedblright{}}}
\noindent \textquotedblleft{\textit{Um bilhete?\textquotedblright{}}}
\noindent \textquotedblleft{\textit{Dizia: 'Avise a próxima. Ela vai vir atrás de mim.'\textquotedblright{}}}
\noindent Lúcia parou de anotar.
\noindent \textquotedblleft{\textit{A senhora guardou o bilhete?\textquotedblright{}}}
\noindent \textquotedblleft{\textit{Guardei.}\textquotedblright} A idosa abriu uma caixa de costura. No fundo, um papel amarelado, dobrado quatro vezes.
\noindent Lúcia estendeu a mão. O papel tremia na ponta dos dedos dela.
\noindent A caligrafia era a mesma do diário.

\noindent \textbf{A próxima leitora.}
\noindent Ela.
\noindent Ela era a próxima.
\noindent A fila tinha chegado nela.
\noindent Lúcia passou a noite cruzando os dados que coletou. O bilhete de Oliveira, as anotações de Joaquim, o laudo de necropsia, a carta de 1979. Cada documento apontava para a mesma conclusão --- uma conclusão que, como psiquiatra forense, ela se recusava a aceitar.

\noindent Às 2:10, ela abriu três envelopes de cadeia de custódia.

\noindent No primeiro, colocou fotografias ampliadas da letra de Oliveira.
\noindent No segundo, raspou uma amostra mínima da tinta do bilhete com lâmina esterilizada.
\noindent No terceiro, colocou uma fibra solta do pano preto em que o diário fora devolvido.

\noindent Lacrou tudo. Assinou sobre a cola. Fotografou os lacres.

\noindent Às 2:47, enviou cópias criptografadas para uma perita grafotécnica da Polícia Civil e para um laboratório particular em Contagem. No assunto, escreveu apenas: \textit{comparação de material histórico --- urgente}.

\noindent Depois guardou o diário de Joaquim em uma caixa de metal, fechou com cadeado e deixou a chave dentro de um copo com água, sobre a mesa.

\noindent Era um gesto absurdo. Mas se acordasse com a chave seca na gaveta, saberia que precisava parar.

\noindent Às 3:14, o celular vibrou.

\noindent A perita respondera rápido demais.

\noindent \textquotedblleft{\textit{A letra do bilhete não é de Oliveira. Também não é falsa. Ligue para mim quando puder. Não leia mais nada até falarmos.\textquotedblright}}

\noindent Lúcia olhou para a caixa de metal. A água do copo estava morna. A chave não estava lá.

\vspace{0.5em}
\noindent\textit{\textbf{↪ Rotas:} \rota{DOC-01} \rotasep \rota{LUC-02} \rotasep \rota{EPI-01}}
